\section{Background}
\label{sec:background}

\draftnote{The physiological background (PLAN.md task 1.7, \file{docs/background.md}) is the scientific lead's
section and has not been drafted yet. The text below is a placeholder written from the corpus and the design
document so that the report reads end to end; it should be replaced or rewritten by the co-author.}

\subsection{Unloading atrophy is not uniform across muscles}

Bed-rest studies that image individual muscles find that loss is concentrated in the muscles that carry body
weight in upright stance. In the first Berlin Bed-Rest Study, MRI of the whole lower limb showed the calf, and the
soleus and gastrocnemius in particular, among the most affected muscles, with the vasti losing more than rectus
femoris and the hip musculature losing less~\citep{belavy2009}. The same group showed that atrophy is heterogeneous
even within a single muscle, depending on where along its length the slice is taken~\citep{miokovic2012}. The
90-day Long-Term Bed Rest study at MEDES documented large losses in the knee extensors and plantar flexors and their
partial prevention by flywheel resistance exercise~\citep{alkner2004,belavy2017}, and a 17-week horizontal bed rest
at NASA established the longest unloading duration in the corpus~\citep{leblanc1992}. A systematic review of
bed-rest studies described the loss of strength and muscle size as nonuniform across muscle groups and over
time~\citep{marusic2021}.

Two explanations are usually offered for this selectivity: fibre-type composition, since the soleus is dominated
by slow (type~I) fibres, and mechanical role, since the soleus and the vasti resist collapse of the ankle and knee
under body weight. The two are confounded in most muscles, but not in all. Tibialis anterior is also slow-fibre
dominant~\citep{johnson1973} yet lifts the foot rather than carrying body weight, and rectus femoris extends the
knee but also flexes the hip. How the framework classifies these muscles, and why, is set out in
\cref{sec:muscle-map}.

\subsection{Countermeasures}

Most recent campaigns include one or more countermeasure arms: resistance or flywheel exercise, aerobic exercise,
reactive jumps, whole-body vibration, artificial gravity by short-arm centrifugation, neuromuscular electrical
stimulation, blood-flow restriction and nutritional supplementation, alone or in combination. The corpus contains
examples of each~\citep[e.g.][]{zange2009,kramer2017,mandic2026,hansen2024,fuchs2025bfr,trappe2024sprint}. At the
number of campaigns available here, countermeasures can be modelled only coarsely (\cref{sec:features}).

\subsection{Pooling aggregate data from clustered studies}

The dataset consists of published group means, several per paper and several papers per campaign. The standard
statistical treatment of such data is aggregate-data meta-regression with a multilevel random-effects structure:
effect sizes nested in studies, studies nested in higher-level clusters~\citep{assink2016,harrer2021,viechtbauer2010}.
Its practical advantage here is that it does not require the within-study covariances between muscles, which no
bed-rest paper reports. Dose--response meta-analysis offers flexible, spline-based forms for the exposure
term~\citep{crippa2019}. Two cautions from that literature shape the design: the number of study-level covariates
that can be supported is roughly one per ten independent studies~\citep{higgins2011}, and many published
meta-regressions fail on exactly the points of clustering, covariate count and ecological
inference~\citep{geissbuhler2021}.

\subsection{Machine learning at small sample sizes}

Flexible learners such as random forests~\citep{breiman2001}, gradient boosting~\citep{friedman2001} and support
vector regression~\citep{smola2004} are data hungry: simulation studies suggest that they need many times more
events per variable than classical regression before their performance stabilises~\citep{vanderploeg2014}. When
observations are clustered, random cross-validation splits place correlated observations on both sides of the fold
boundary and overstate performance; grouped (blocked) cross-validation is the recommended
remedy~\citep{roberts2017}. Tabular foundation models take a different route to small data: TabPFN is a transformer
pretrained on millions of synthetic datasets that predicts for a new table in context, and was designed for tables
of up to 10,000 rows~\citep{hollmann2025}. Model-agnostic importance measures such as SHAP~\citep{lundberg2017} describe what a
fitted model relies on, not what causes the outcome, and at small sample sizes their rankings can be unstable
across folds.

\subsection{Positioning: what machine learning adds to meta-regression}

The sharpest methodological question this project faces is why machine learning should be used at all when
mixed-effects meta-regression is the established tool. The position taken here is that it should not replace the
meta-regression, and does not. The meta-regression (tier~1) is the estimator: it produces coefficients with
intervals, and it is the primary result. The machine-learning families (tier~2) are a test: given the same rows
and the same leave-one-campaign-out discipline, do flexible models find structure---interactions, thresholds,
nonlinearities in covariates---that the curve misses? A null answer is informative, because it says how much
signal the published literature contains. The language-model forecast (tier~3) asks a different question again:
whether a model that reads the \emph{description} of a campaign, and the other campaigns' data in context, can
extract information the numeric columns do not carry. Its outputs are probability distributions and are scored with
proper scoring rules~\citep{gneiting2007,epstein1969}.
